\section{Detailed numerical results and auxiliary decisions}
\label{app:detailed-results}

The audit and extension tables are generated from validated summary JSON artifacts,
including negative seeds and failed searches; exact paths and hashes are
recorded in Appendix~\ref{app:experimental-details}.

\subsection{A non-permutation chart is visible in optimizer response}

% Auto-generated by experiments/generate_paper_audit_tables.py; do not edit.
% Source: results/response_summary.json
% Source SHA256: fd3f8ade6b43e6830e5a5c6bdc188509b77da844c4fa977c5d25e5af1e3d3224
\providecommand{\ResponseProbeDetected}{24}
\providecommand{\ResponseProbeExpected}{24}
\providecommand{\ResponseEffectiveMax}{\ensuremath{1.70\!\times\!10^{-16}}}
\providecommand{\ResponseFormulaMax}{\ensuremath{9.57\!\times\!10^{-13}}}
\providecommand{\ResponsePermutationMax}{\ensuremath{8.82\!\times\!10^{-17}}}
\providecommand{\ResponseMedianMin}{0.5045}
\providecommand{\ResponseMedianMax}{0.5058}
\begin{table}[H]
\centering
\scriptsize
\caption{Frozen-checkpoint response audit.  The effective-product column checks the fixed non-permutation gauge; $\Delta\mathcal R$ compares Euclidean instantaneous response operators on fixed Gaussian effective-gradient probes.  The permutation column is the label-swap control.}
\label{tab:response-audit}
\resizebox{\linewidth}{!}{%
\begin{tabular}{crrrrrrr}
\toprule
Seed & Effective product rel. error & $\min A$ & $\sigma_{\min}(A)$ & $\sigma_{\min}(B)$ & Perm. max rel. $\Delta\mathcal R$ & Nonperm. median rel. $\Delta\mathcal R$ & Detected probes \\
\midrule
0 & $1.58\!\times\!10^{-16}$ & $2.88\!\times\!10^{-4}$ & 1.7291 & 27.445 & $8.82\!\times\!10^{-17}$ & 0.5045 & 8/8 \\
1 & $1.70\!\times\!10^{-16}$ & $2.72\!\times\!10^{-4}$ & 1.7291 & 27.339 & $8.82\!\times\!10^{-17}$ & 0.5051 & 8/8 \\
2 & $1.63\!\times\!10^{-16}$ & $2.61\!\times\!10^{-4}$ & 1.7290 & 27.371 & $8.82\!\times\!10^{-17}$ & 0.5058 & 8/8 \\
\bottomrule
\end{tabular}%
}
\end{table}


Table~\ref{tab:response-audit} checks the response formula and separates
a fixed alternative at frozen checkpoints under
Theorem~\ref{thm:response_stabilizer}'s assumptions.  All stated rank conditions hold, the
non-permutation reparameterization preserves the effective model to at most
\ResponseEffectiveMax{} relative error.  On the declared $4\times32$
extraction, the analytic response agrees with autograd within numerical
precision, and the worst finite-difference error across native and transformed
charts is \ResponseFormulaMax{}.  The maximum permutation-control discrepancy
is \ResponsePermutationMax{} after rounding, whereas all
\ResponseProbeDetected{}/\ResponseProbeExpected{} checkpoint--probe
evaluations detect a change (eight per checkpoint, with overlapping seed
windows and 10 unique Gaussian draws; per-seed medians
$[\ResponseMedianMin{},\ResponseMedianMax{}]$).
The separate structured audit reduces the all-gradient quantifier to
\ResponseTomographyProbeCount{} designed analytic response evaluations, which
reconstruct all response components with worst relative error
\ResponseTomographyReconstructionMax{}; the
permutation control stays below \ResponseTomographyPermutationMax{}; the
minimum across seeds of each seed's maximum fixed non-permutation difference
over the four designed probes is \ResponseTomographyNonpermutationMin{} after
rounding.  This does not say that every individual probe reaches that value.
These are analytic evaluations of
Equation~\ref{eq:response_blocks}, validated on the extraction, not four
full-model JVP calls.  The eight Gaussian probes remain a fixed-alternative
test; the designed certificate assumes exact common product and the stated
response family.  Neither audit describes AdamW or a historical graph.

\paragraph{Independent full-dimensional response measurements.}
A subsequent audit obtains the responses by differentiating the complete
packed map $(Z,B)\mapsto\operatorname{softmax}(Z/\tau)B$ at all three
released checkpoints. Reverse-mode differentiation of a linear probe loss
gives the factor gradients; a forward-mode JVP then measures the induced
effective-parameter response. The observation code does not call the analytic
response formula. Each native, permutation, and fixed non-permutation chart
uses four designed probes and two unseen Gaussian probes.
Table~\ref{tab:autograd-tomography} reports formula agreement and prediction of
the unseen responses from the reconstructed operator. The maximum relative
formula discrepancy is $3.71\times10^{-16}$, and the maximum held-out response
error is $5.08\times10^{-14}$. A five-level observation-noise grid from zero
to $10^{-4}$ relative Frobenius norm satisfies the deterministic component
error bounds, with recorded floating-point slack. These measurements concern
the full shared-parameter map, not an end-to-end language-model loss or an
AdamW step; operator reconstruction does not itself implement noisy factor
recovery. The additive records preserve the original analytic audit.



\paragraph{Exploratory reconstruction of the factors.}
The construction in Appendix~\ref{app:response_tomography} recovers factors
from the effective product and measured components, without supplying the
original factors to the reconstruction. We test it on the native charts of
the same three checkpoints, using fresh autodiff observations and one fixed
Gaussian weight vector for joint diagonalization. Table~\ref{tab:factor-recovery}
reports the errors after a common permutation chosen only for evaluation.
The noiseless combined relative errors are below $3.7\times10^{-9}$.
All five noise levels complete without a rank or eigengap failure, but the
unprojected noisy estimates need not have exact unit row sums. These are
three checkpoint diagnostics with exact effective products, not independent
training replications, adversarial-noise tests, or a general stability result
for the PSD-clipped recovery algorithm.

% Generated by experiments/summarize_factor_recovery.py.
\begin{table}[H]
\centering\small
\caption{Exploratory factor recovery from full-dimensional autodiff observations. Orbit error combines relative router and basis errors after one common permutation. Noisy rows use relative observation noise $10^{-4}$ and no simplex projection.}
\label{tab:factor-recovery}
\begin{tabular}{@{}crrr@{}}
\toprule
Seed & Noiseless orbit error & Noisy orbit error & Noisy row-sum error \\
\midrule
0 & 3.66e-09 & 1.81e-05 & 6.06e-05 \\
1 & 3.23e-09 & 1.61e-05 & 6.68e-05 \\
2 & 3.28e-09 & 1.59e-05 & 5.99e-05 \\
\bottomrule
\end{tabular}
\end{table}


\subsection{Gauge-equivalent ASLoRA states change the first merge decision}
\label{sec:aslora-results}

% Auto-generated by experiments/generate_paper_audit_tables.py; do not edit.
% Source: results/aslora_summary.json
% Source SHA256: 1be80d31ca4208c414698c3b5fbd45bfb9d51835aa8f16f8cb43bab082840c41
\providecommand{\ASLORAAdjacentPrimaryChanges}{11}
\providecommand{\ASLORAAllPrimaryChanges}{14}
\providecommand{\ASLORATotalGaugeInstances}{30}
\providecommand{\ASLORAGaugesPerSeed}{10}
\begin{table}[H]
\centering
\small
\setlength{\tabcolsep}{3pt}
\caption{ASLoRA first-merge decision audit over \ASLORAGaugesPerSeed{} gauges per seed.  Changes are relative to the identity-gauge raw-running-$B$ action.  The loss interval is the canonical validation difference (selected action minus identity action), not a gauged float32 re-execution; accuracy and F1 intervals are in Appendix Table~\ref{tab:aslora-metric-details}.  The chain column requires a changed action with a nonzero validation consequence in every seed.}
\label{tab:aslora-chain}
\begin{tabular}{@{}llcrrcc@{}}
\toprule
Candidates & Scope & By seed & Changed & Seeds & $\Delta$ loss & Chain \\
\midrule
Adjacent & per-proj. & $(2,6,3)$ & 11/30 & 3/3 & $[-0.0531,+0.0089]$ & yes \\
Adjacent & joint pair & $(0,5,0)$ & 5/30 & 1/3 & $[+0.0000,+0.0077]$ & no \\
All pairs & per-proj. & $(4,7,3)$ & 14/30 & 3/3 & $[-0.0531,+0.0027]$ & yes \\
All pairs & joint pair & $(0,0,0)$ & 0/30 & 0/3 & $[+0.0000,+0.0000]$ & no \\
\bottomrule
\end{tabular}
\end{table}


The primary per-projection chain replicates under both readings of the source's
candidate set.  Among \ASLORATotalGaugeInstances{} predeclared gauge instances,
the raw-running-$B$ action changes in
\ASLORAAdjacentPrimaryChanges{} instances for adjacent candidates and
\ASLORAAllPrimaryChanges{} instances for all pairs.  Every seed has at least
one changed action with a nonzero canonical validation consequence under both
candidate readings.  The signs in Table~\ref{tab:aslora-chain} go in both
directions: the experiment demonstrates representation-dependent decisions,
not that every gauge-induced change is harmful or that the identity action is
optimal.  The joint-same-pair sensitivity does not replicate (one seed for
adjacent candidates and none for all pairs), and Appendix
Table~\ref{tab:aslora-baselines} shows that neither gauge-invariant rule
uniformly improves on raw running \(B\) across all seeds under both candidate
readings.  The causal contrast uses a
float64-equivalent snapshot only to select layer indices, then executes the
same temporary tie in the unchanged canonical model; factorized-float32
re-execution is a passing secondary diagnostic.  Thus the complete chain is
concrete but local to one first merge, method, task, and fixed gauge grid.

\subsection{Prospective controls: learned routers and recovery training}
\label{sec:enhanced-decisions}

The historical witness in Appendix~\ref{app:historical-fold-witness}
motivates a separately registered extension;
it does not supply its outcomes. We train three fresh width-128, twelve-layer,
four-basis byte LMs for a fixed 6,000 steps with random router initialization
and a separate router learning rate. Router L1 movement is $0.593$--$0.729$;
all three models pass the unigram-improvement and router-movement gates.
Both gauge families receive 5,000 router-only trials, now requiring the same
sorted group-size multiset as the native partition. Positive-stochastic
search finds a first valid candidate in every seed; signed-local search fails
in all three. Failures are retained and never replaced by larger searches.

\begin{table}[t]
\centering\small
\caption{Prospective equal-size folding audit. Signed BPB differences are gauge-selected minus native-selected folds after matched recovery updates. Each row is one independent training seed. Positive-stochastic search succeeds in 3/3 seeds; signed-local search fails in 3/3, each with 5,000 router-only candidates.}
\label{tab:enhanced-decisions}
\begin{tabular}{rrrrr}
\toprule
Seed & Router movement & $\Delta_0$ & $\Delta_{100}$ & $\Delta_{500}$ \\
\midrule
10 & 0.593 & -0.00112 & +0.00224 & +0.00009 \\
11 & 0.720 & +0.87112 & +0.03574 & +0.01823 \\
12 & 0.729 & +0.14873 & +0.00580 & +0.00343 \\
\bottomrule
\end{tabular}
\end{table}


Every fold uses the same native effective tensors. Hard routes remain fixed
while all remaining parameters receive identical minibatches, optimizer
initialization, update counts, and learning rates. Table~\ref{tab:enhanced-decisions}
shows substantial immediate differences in two seeds, but much smaller
differences after 500 recovery updates. Conditional paired 95\% bootstrap
intervals at that stage are $[-0.00094,0.00112]$, $[0.01654,0.01989]$, and
$[0.00244,0.00437]$ BPB, respectively; they describe the fixed evaluation
blocks, not uncertainty over training populations. The same-size
effective-parameter local-search baseline selects the native partition in
two seeds and performs slightly worse after recovery in the third.
Thus removing the earlier initialization and size confounds strengthens the
decision-sensitivity evidence, without establishing persistent large harm or
general superiority of an invariant rule. This is one architecture and corpus,
with a fixed training budget rather than a convergence certificate.

\subsection{Natural task gradients and finite arithmetic}
\label{sec:natural-gradients}

We differentiate actual byte-LM cross-entropy with respect to independent
layer-effective weights, holding other parameters fixed. A separate forward
and chain-rule check agrees with the original model. At each of the three
original and three new checkpoints, 16 fixed training minibatches supply
natural directions. Controlled reverse/forward autodiff measures the response;
the first $q\in\{1,4,8,12\}$ directions fit the estimator, and the last four
test prediction. The estimator uses the exact product row space and known
response structure, not the true factors.



At $q=4$, all six checkpoints meet the rank conditions of
Proposition~\ref{prop:natural-probes}. The largest held-out response error is
below $10^{-13}$ and the factor error below $3.1\times10^{-8}$.
This is stronger than separation of a single alternative chart, but remains
conditional on the known Euclidean response family and precise measurements.
A single direction has insufficient local rank; at the new checkpoints its
factor errors range from $0.258$ to $1.34$. Generic linear prediction from the
observed gradient span also differs from the structural estimator. These
directions are measurements at fixed checkpoints, not observations of an
uncontrolled training trajectory.

Explicit finite Euclidean steps expose a numerical boundary. At step size
$10^{-5}$, the new checkpoints' FP64 response errors are
$3.3\times10^{-8}$--$5.2\times10^{-8}$, whereas FP32 subtraction errors are
$0.0208$--$0.0260$ and BF16 errors are about one.
Here factors and arithmetic are cast to the stated precision; this is not a
test of BF16 training with FP32 master weights. Appendix
Figure~\ref{fig:enhancement-precision} reports the step-size tradeoff and the
separate joint-noise conditioning diagnostic.

\subsection{Noisy recovery and the limits of observable rejection}
\label{sec:trusted-recovery}

We fix a prospective development/calibration/evaluation split before
measuring the final outcomes. Forty held-out synthetic seeds each contribute
48 correlated noise/probe/regime variants; four freshly trained byte-LM
checkpoints each contribute 12. Both product and autodiff response
observations receive relative noise from zero to $10^{-2}$, and the rank-$K$
product row subspace is re-estimated from the noisy product.
Three fixed methods use the same observations: spectral recovery,
simplex projection with a basis refit, and three-start constrained joint
least squares. The estimator never receives the true factors.
Appendix~\ref{app:trusted-experiments} gives splits, budgets, calibration,
and all diagnostic comparators.

\begin{table}[t]
\centering
\small
\caption{Frozen held-out noisy recovery. Each synthetic instance contributes 48 correlated variants; each language checkpoint contributes 12. Acceptance uses the calibrated empirical score. Bad means factor error $>0.10$; good means $\le0.05$. Local passes concern only a numerical restricted-subspace component.}
\label{tab:trusted-recovery}
\begin{tabular}{llrrrrr}
\toprule
Domain & Method & Median error & Bad/total & Accept & Bad/accept & Local \\
\midrule
Synthetic & Spectral & 0.0041 & 564/1920 & 196 & 0/196 & 0 \\
Synthetic & Projected & 0.0033 & 566/1920 & 1396 & 56/1396 & 0 \\
Synthetic & Joint fit & 0.0001 & 227/1920 & 1782 & 92/1782 & 140 \\
Byte-LM & Spectral & 0.00017 & 1/48 & 0 & -- & 0 \\
Byte-LM & Projected & 0.0002 & 1/48 & 48 & 1/48 & 0 \\
Byte-LM & Joint fit & 3.3e-05 & 0/48 & 48 & 0/48 & 2 \\
\bottomrule
\end{tabular}
\end{table}


Joint fitting reduces the synthetic median factor error from
$3.34\times10^{-3}$ for projected spectral recovery to $1.01\times10^{-4}$,
and the number of errors above $0.10$ from $566$ to $227$ out of $1920$.
At the four byte-LM checkpoints, all 48 joint-fit errors are below $0.0107$,
including the $1\%$ observation-noise cases. These are controlled Euclidean
measurements at checkpoints, not noisy AdamW trajectory observations.

The more demanding question is whether an observable diagnostic can tell
when to trust a result. With its threshold fixed on separate calibration
data, the combined residual/conditioning/subspace/multi-start score accepts
$1782/1920$ synthetic outputs, including 92 bad outputs. Its $5.16\%$
accepted-bad fraction has a seed-cluster bootstrap interval of
$[4.50\%,5.79\%]$ and has a point estimate above the $5\%$ calibration target;
the interval includes that target.
Residual-only rejection accepts $1757/1920$ with $68/1757=3.87\%$ bad.
Thus, the combined score does not establish superiority.
The full evaluation retains a joint-fit error of $10.54$, and a falsely
accepted error of $0.82$. These failures prevent a general reliability claim.

The explicit local separation bound in
Proposition~\ref{prop:trusted-local} addresses a different question.
Only $140/1920$ synthetic and $2/48$ language outputs pass its numerical
local-component conditions. These passes neither establish truth membership
in the restricted subspace nor rule out distant solutions.
